\begin{problem}[第35届全国中学生物理竞赛决赛][弦的受迫振动与驻波]{99}
    如图，一张紧的弦沿 $x$ 轴水平放置，长度为 $L$ 。弦的左端位于坐标原点。弦可通过其左、右端与振源连接，使弦产生沿 $y$ 方向的横向受迫振动，振动传播的速度为 $u$ 。
    \begin{subproblem}
        固定弦的右端 $P_{2}$ ，将其左端 $P_{1}$ 与振源连接，稳定时，左端 $P_{1}$ 的振动表达式为 $y(x = 0,t) = A_0\cos (\omega t)$ ，其中 $A_0$ 为振幅， $\omega$ 为圆频率。
        \begin{subsubproblem}
            已知弦上横波的振幅在传播方向上有衰减，衰减常量为 $\gamma (\gamma >0)$ ，求弦上各处振动的振幅；（已知：在无限长弦上沿 $x$ 轴正方向传播的振幅逐渐衰减的横波表达式为 $y(x,t) = Ae^{-\gamma x}\cos \left(\omega t - \frac{\omega x}{u} +\varphi\right)$ ，其中 $A$ 和 $\varphi$ 分别为 $x = 0$ 处振动的振幅和初相位。）
        \end{subsubproblem}
        \begin{subsubproblem}
            忽略波的振幅在传播方向上的衰减，求弦上驻波的表达式，并确定其波腹和波节处的 $x$ 坐标。
        \end{subsubproblem}
    \end{subproblem}
    \begin{subproblem}
        将 $P_{1}$ 、 $P_{2}$ 都与振源连接， $P_{1}$ 、 $P_{2}$ 处的振动表达式分别为： $y(x = 0,t) = A_0\cos \omega t$ 、 $y(x = L,t) = A_0\cos (\omega t + \varphi_0)$ ，其中 $\varphi_0$ 为常量。忽略波的振幅在传播方向上的衰减，分别计算 $\varphi_0 = 0$ 和 $\varphi_0 = \pi$ 情形下弦上各处振动的表达式以及共振时圆频率 $\omega$ 应满足的条件。
    \end{subproblem}
\end{problem}